\textbf{Phần 1. Trắc nghiệm lựa chọn}
\setcounterpageref{ex}{0}

\begin{ex} %[5H9N52]
Muốn tính thể tích hình hộp chữ nhật có chiều dài $a$, chiều rộng $b$ và chiều cao $c$ (cùng đơn vị đo), ta dùng công thức
\choice
{\True $V = a \times b \times c$}
{$V = (a + b) \times c$}
{$V = (a + b) \times 2 \times c$}
{$V = a + b + c$}
\loigiai{
$V = a \times b \times c$.\\
\textbf{Đáp án của câu hỏi là A}
}
\end{ex}

\begin{ex} %[5H9N52]
Một hình hộp chữ nhật có chiều dài $7\text{ cm}$, chiều rộng $4\text{ cm}$, chiều cao $5\text{ cm}$. Thể tích của hình hộp đó là
\choice
{$68\text{ cm}^3$}
{\True $140\text{ cm}^3$}
{$140\text{ cm}^2$}
{$120\text{ cm}^3$}
\loigiai{
$V = 7 \times 4 \times 5 = 140\text{ (cm}^3\text{)}$.\\
\textbf{Đáp án của câu hỏi là B}
}
\end{ex}

\begin{ex} %[5H9H52]
Một bể nước hình hộp chữ nhật có diện tích đáy là $2{,}4\text{ m}^2$, chiều cao $1{,}5\text{ m}$. Thể tích của bể nước là
\choice
{\True $3{,}6\text{ m}^3$}
{$3{,}6\text{ m}^2$}
{$36\text{ m}^3$}
{$360\text{ l}$}
\loigiai{
$V = S_{đáy} \times h = 2{,}4 \times 1{,}5 = 3{,}6\text{ (m}^3\text{)}$.\\
\textbf{Đáp án của câu hỏi là A}
}
\end{ex}

\begin{ex} %[5H9H52]
Một hình hộp chữ nhật có thể tích $180\text{ cm}^3$, chiều dài $9\text{ cm}$, chiều rộng $4\text{ cm}$. Chiều cao của hình hộp chữ nhật là
\choice
{$4\text{ cm}$}
{$6\text{ cm}$}
{$8\text{ cm}$}
{\True $5\text{ cm}$}
\loigiai{
Chiều cao: $180 : (9 \times 4) = 180 : 36 = 5\text{ (cm)}$.\\
\textbf{Đáp án của câu hỏi là D}
}
\end{ex}

\begin{ex} %[5H9H52]
Một khối gỗ hình hộp chữ nhật có kích thước dài $1{,}2\text{ m}$, rộng $0{,}5\text{ m}$, cao $0{,}4\text{ m}$. Thể tích khối gỗ đó là
\choice
{\True $240\text{ dm}^3$}
{$0{,}24\text{ dm}^3$}
{$24\text{ dm}^3$}
{$2\,400\text{ dm}^3$}
\loigiai{
$V = 1{,}2 \times 0{,}5 \times 0{,}4 = 0{,}24\text{ m}^3 = 240\text{ dm}^3$.\\
\textbf{Đáp án của câu hỏi là A}
}
\end{ex}

\begin{ex} %[5H9H52]
Một bể cá hình hộp chữ nhật có chiều dài $60\text{ cm}$, chiều rộng $40\text{ cm}$ và chiều cao $50\text{ cm}$. Dung tích tối đa của bể cá là
\choice
{$120\text{ l}$}
{$12\text{ l}$}
{\True $120\text{ l}$}
{$1\,200\text{ l}$}
\loigiai{
$V = 6 \times 4 \times 5 = 120\text{ dm}^3 = 120\text{ l}$.\\
\textbf{Đáp án của câu hỏi là C}
}
\end{ex}

\begin{ex} %[5H9V52]
Một bể nước hình hộp chữ nhật có chiều dài $2\text{ m}$, chiều rộng $1{,}2\text{ m}$, chiều cao $1{,}5\text{ m}$. Bể đang chứa nước đến $60\%$ chiều cao bể. Lượng nước hiện có trong bể là
\choice
{$3{,}6\text{ m}^3$}
{$2\,400\text{ l}$}
{$1{,}8\text{ m}^3$}
{\True $2{,}16\text{ m}^3$}
\loigiai{
Thể tích bể: $2 \times 1{,}2 \times 1{,}5 = 3{,}6\text{ (m}^3\text{)}$.\\
Lượng nước hiện có: $3{,}6 \times 60\% = 2{,}16\text{ (m}^3\text{)}$.\\
\textbf{Đáp án của câu hỏi là D}
}
\end{ex}

\begin{ex} %[5H9V52]
Một cái hố hình hộp chữ nhật sâu $1{,}5\text{ m}$, đáy là hình vuông có chu vi $8\text{ m}$. Thể tích đất đã đào lên từ cái hố đó là
\choice
{$12\text{ m}^3$}
{\True $6\text{ m}^3$}
{$8\text{ m}^3$}
{$4\text{ m}^3$}
\loigiai{
Cạnh đáy: $8 : 4 = 2\text{ (m)}$.\\
Diện tích đáy: $2 \times 2 = 4\text{ (m}^2\text{)}$.\\
Thể tích đất: $4 \times 1{,}5 = 6\text{ (m}^3\text{)}$.\\
\textbf{Đáp án của câu hỏi là B}
}
\end{ex}

\begin{ex} %[5H9V52]
Người ta đổ vào một bể nước rỗng hình hộp chữ nhật $1\,800\text{ l}$ nước thì mực nước dâng cao $0{,}75\text{ m}$. Diện tích đáy của bể nước là
\choice
{$2\text{ m}^2$}
{$2{,}5\text{ m}^2$}
{\True $2{,}4\text{ m}^2$}
{$3\text{ m}^2$}
\loigiai{
$1\,800\text{ l} = 1{,}8\text{ m}^3$.\\
Diện tích đáy: $1{,}8 : 0{,}75 = 2{,}4\text{ (m}^2\text{)}$.\\
\textbf{Đáp án của câu hỏi là C}
}
\end{ex}

\begin{ex} %[5H9V52]
Nếu chiều dài tăng gấp $2$ lần, chiều rộng tăng gấp $3$ lần và chiều cao không đổi thì thể tích của hình hộp chữ nhật tăng lên gấp mấy lần?
\choice
{$5$ lần}
{\True $6$ lần}
{$8$ lần}
{$9$ lần}
\loigiai{
Thể tích mới: $V' = (2 \times a) \times (3 \times b) \times c = 6 \times (a \times b \times c) = 6V$.\\
\textbf{Đáp án của câu hỏi là B}
}
\end{ex}

\begin{ex} %[5H9V52]
Một bể kính nuôi cá dạng hình hộp chữ nhật có kích thước lòng trong: dài $80\text{ cm}$, rộng $50\text{ cm}$, cao $45\text{ cm}$. Mực nước ban đầu cao $30\text{ cm}$. Sau khi thả một hòn non bộ vào thì mực nước dâng lên cao $38\text{ cm}$ (hòn non bộ ngập hoàn toàn). Thể tích của hòn non bộ là
\choice
{$16\text{ dm}^3$}
{$24\text{ dm}^3$}
{$40\text{ dm}^3$}
{\True $32\text{ dm}^3$}
\loigiai{
Chiều cao nước dâng thêm: $38 - 30 = 8\text{ (cm)}$.\\
Thể tích hòn non bộ: $80 \times 50 \times 8 = 32\,000\text{ (cm}^3\text{)} = 32\text{ dm}^3$.\\
\textbf{Đáp án của câu hỏi là D}
}
\end{ex}

\begin{ex} %[5H9C52]
Một bể bơi hình hộp chữ nhật dài $20\text{ m}$, rộng $12\text{ m}$, sâu $1{,}6\text{ m}$. Người ta dùng $2$ máy bơm cùng bơm nước vào bể, mỗi phút máy 1 bơm được $120\text{ l}$, máy 2 bơm được $180\text{ l}$. Hỏi sau bao lâu cả hai máy bơm đầy được $75\%$ bể nước?
\choice
{$12\text{ giờ}$}
{$18\text{ giờ}$}
{\True $16\text{ giờ}$}
{$20\text{ giờ}$}
\loigiai{
Thể tích bể: $20 \times 12 \times 1{,}6 = 384\text{ m}^3 = 384\,000\text{ l}$.\\
$75\%$ thể tích: $384\,000 \times 75\% = 288\,000\text{ l}$.\\
Mỗi phút cả 2 máy bơm: $120 + 180 = 300\text{ (l)}$.\\
Thời gian bơm: $288\,000 : 300 = 960\text{ (phút)} = 16\text{ giờ}$.\\
\textbf{Đáp án của câu hỏi là C}
}
\end{ex}

\textbf{Phần 2. Trắc nghiệm đúng sai}
\setcounterpageref{ex}{0}

\begin{ex} %[5H9H52]
Một khối gỗ hình hộp chữ nhật có kích thước dài $1{,}5\text{ m}$, rộng $0{,}6\text{ m}$, cao $0{,}4\text{ m}$. Xét tính đúng sai:
\choiceTF
{\True Diện tích mặt đáy của khối gỗ là $0{,}9\text{ m}^2$.}
{\True Thể tích của khối gỗ là $0{,}36\text{ m}^3$.}
{Thể tích khối gỗ tương đương với $36\text{ dm}^3$.}
{\True Thể tích khối gỗ tương đương với $360\text{ dm}^3$.}
\loigiai{
\begin{itemize}
    \item $S_{đáy} = 1{,}5 \times 0{,}6 = 0{,}9\text{ m}^2$. Đúng.
    \item $V = 0{,}9 \times 0{,}4 = 0{,}36\text{ m}^3$. Đúng.
    \item $0{,}36\text{ m}^3 = 360\text{ dm}^3 \neq 36\text{ dm}^3$. Ý 3 Sai, Ý 4 Đúng.
\end{itemize}
\begin{itemize}
    \item a) \textbf{Đúng}.
    \item b) \textbf{Đúng}.
    \item c) \textbf{Sai}.
    \item d) \textbf{Đúng}.
\end{itemize}
}
\end{ex}

\begin{ex} %[5H9V52]
Một bể nước hình hộp chữ nhật có chiều dài $2\text{ m}$, chiều rộng $1{,}5\text{ m}$ và chiều cao $1{,}2\text{ m}$. Bể đang chứa $1\,800\text{ l}$ nước. Xét tính đúng sai:
\choiceTF
{\True Dung tích toàn bộ của bể là $3\,600\text{ l}$.}
{\True Mực nước hiện tại trong bể cao $0{,}6\text{ m}$.}
{Lượng nước hiện tại chiếm $75\%$ thể tích bể.}
{\True Cần đổ thêm $1\,800\text{ l}$ nước nữa thì đầy bể.}
\loigiai{
\begin{itemize}
    \item Dung tích: $2 \times 1{,}5 \times 1{,}2 = 3{,}6\text{ m}^3 = 3600\text{ l}$. Đúng.
    \item Mực nước: $1{,}8 : (2 \times 1{,}5) = 1{,}8 : 3 = 0{,}6\text{ m}$. Đúng.
    \item Tỉ lệ nước hiện có: $1800 : 3600 = 50\% \neq 75\%$. Sai.
    \item Cần thêm: $3600 - 1800 = 1800\text{ l}$. Đúng.
\end{itemize}
\begin{itemize}
    \item a) \textbf{Đúng}.
    \item b) \textbf{Đúng}.
    \item c) \textbf{Sai}.
    \item d) \textbf{Đúng}.
\end{itemize}
}
\end{ex}

\textbf{Phần 3. Trả lời ngắn}
\setcounterpageref{ex}{0}

\begin{ex} %[5H9H52]
Một hình hộp chữ nhật có chiều dài $12\text{ cm}$, chiều rộng $8\text{ cm}$, chiều cao $5\text{ cm}$. Tính thể tích theo xăng-ti-mét khối.
\shortans{480}
\loigiai{
\textbf{Đáp số:} $480$\\
$12 \times 8 \times 5 = 480\text{ cm}^3$.
}
\end{ex}

\begin{ex} %[5H9H52]
Một khối bê tông hình hộp chữ nhật có diện tích đáy là $3{,}5\text{ m}^2$, chiều cao $0{,}8\text{ m}$. Tính thể tích khối bê tông theo mét khối.
\shortans{2,8}
\loigiai{
\textbf{Đáp số:} $2{,}8$\\
$3{,}5 \times 0{,}8 = 2{,}8\text{ m}^3$.
}
\end{ex}

\begin{ex} %[5H9V52]
Một bể nước hình hộp chữ nhật dài $1{,}8\text{ m}$, rộng $1{,}5\text{ m}$, sâu $1{,}2\text{ m}$. Hỏi bể chứa được tối đa bao nhiêu lít nước?
\shortans{3240}
\loigiai{
\textbf{Đáp số:} $3240$\\
$1{,}8 \times 1{,}5 \times 1{,}2 = 3{,}24\text{ m}^3 = 3\,240\text{ l}$.
}
\end{ex}

\begin{ex} %[5H9V52]
Một hình hộp chữ nhật có thể tích $240\text{ cm}^3$, chiều dài $8\text{ cm}$, chiều cao $6\text{ cm}$. Tính chiều rộng của hình hộp chữ nhật theo xăng-ti-mét.
\shortans{5}
\loigiai{
\textbf{Đáp số:} $5$\\
$240 : (8 \times 6) = 5\text{ cm}$.
}
\end{ex}

\begin{ex} %[5H9V52]
Một bể cá có diện tích đáy là $1\,500\text{ cm}^2$. Người ta đổ vào bể $45\text{ l}$ nước. Tính chiều cao mực nước trong bể theo xăng-ti-mét.
\shortans{30}
\loigiai{
\textbf{Đáp số:} $30$\\
$45\text{ l} = 45\,000\text{ cm}^3 \Rightarrow 45\,000 : 1\,500 = 30\text{ cm}$.
}
\end{ex}

\begin{ex} %[5H9C52]
Một khối kim loại hình hộp chữ nhật có chiều dài $20\text{ cm}$, chiều rộng $15\text{ cm}$, chiều cao $10\text{ cm}$. Biết $1\text{ dm}^3$ kim loại đó nặng $7{,}8\text{ kg}$. Tính khối lượng khối kim loại theo ki-lô-gam.
\shortans{23,4}
\loigiai{
\textbf{Đáp số:} $23{,}4$\\
Thể tích: $20 \times 15 \times 10 = 3\,000\text{ cm}^3 = 3\text{ dm}^3$.\\
Khối lượng: $3 \times 7{,}8 = 23{,}4\text{ kg}$.
}
\end{ex}
